\documentclass[11pt]{article}
\usepackage[margin=1in]{geometry}
\usepackage{amsmath,amssymb,xurl}
\newcommand{\Tr}{\operatorname{Tr}}
\title{Formation scales, the coordinate-rest clock and a light-radiation capacity bound}
\author{}
\date{}
\begin{document}
\maketitle

\paragraph{Scope [R].}
This note serves the formation-clock and matched-radius expansions, the
coordinate-rest free-fall comparator, and the light-field formation budget
named in task six(a). The first two calculations concern a smooth
continuum family and its exterior, respectively. The third bounds a
specified outgoing band and arrival interval; it does not bound every
pair created inside the object. No numerical evolution, coefficient fit
or scientific code is part of this route.

\paragraph{A smooth birth has a quadratic width correction.}
The averaging width should be varied without confusing a displaced birth
with a change of the radius at a fixed clock. Fix a reference radius
$r_s$ and a common boundary clock, and put
\[
 x=r/r_s,\qquad \tau=t/r_s,\qquad
 \varepsilon=\sigma^2/r_s^2,\qquad F=N+v
\]
on the inward-flow branch. The chosen root is that of the specified
dense or applied field throughout; switching the two is not a width
comparison. Let the normalized averaging kernel be even and
cubic-symmetric, with per-axis variance $\sigma^2$ and a uniform
absolute fourth-moment bound
\[
 \int |y|^4|K_\sigma(y)|\,d^3y=O(\sigma^4).
\]
A fixed-shape family $K_\sigma(y)=\sigma^{-3}K(y/\sigma)$ with a
finite absolute fourth moment suffices. On a smooth bare field with
bounded fourth derivatives, Taylor's formula then gives
\[
 K_\sigma f=f+\tfrac12\sigma^2\Delta f
             +O(\sigma^4\|\nabla^4f\|_\infty).
\]
For the coupled problem the state, constraints and back-averaging also
respond. Assume a regular solution family, with those responses included,
\[
 F(\tau,x;\varepsilon)=F_0(\tau,x)+\varepsilon F_1(\tau,x)
                  +O(\varepsilon^2)
\]
in the derivatives used below, with sufficient uniform differentiability
for the Taylor remainders. The kernel expansion motivates this
parameter; it does not by itself prove differentiability of a nonlinear
solution branch or determine $F_1$. The hydrodynamic limit keeps
$a\ll\sigma$ and $\sigma$ smaller than every geometric gradient length;
using $r_s$ as the unit does not replace this local condition.
The zero-width coefficient is a formal smooth continuum reference,
not a permissible $\sigma=0$ lattice calculation. Uniform use along
$\chi=a/r_s\to0$ additionally requires the corresponding regularity
bounds to be uniform in $\chi$.

Suppose the first selected interior birth is isolated, at $x_0>0$, with
\[
 F_0(\tau_0,x_0)=0,\qquad F_{0,x}(\tau_0,x_0)=0,\qquad
 F_{0,\tau}<0,\quad F_{0,xx}>0.
\]
The two equations $F=F_x=0$ determine the birth. Their Jacobian with
respect to $(\tau,x)$ has determinant $F_{0,\tau}F_{0,xx}\ne0$, so the
birth clock and location have regular expansions
\[
 \tau_b=\tau_0+b_t\varepsilon+O(\varepsilon^2),\qquad
 x_b=x_0+b_r\varepsilon+O(\varepsilon^2),
\]
where, with every derivative evaluated at the unperturbed birth,
\[
 b_t=-\frac{F_1}{F_{0,\tau}},\qquad
 b_r=-\frac{F_{1,x}+F_{0,x\tau}b_t}{F_{0,xx}}.
\]
No universal sign follows. Degenerate contact, boundary birth, a change
of the first competing minimum, a non-smooth source or loss of the
solution branch lies outside this argument.

At a common clock $\tau_m$ where the chosen outer root is simple and
uniformly separated from birth, the corresponding expansion is
\[
 x_t(\tau_m;\varepsilon)
 =x_t(\tau_m;0)
 -\varepsilon\frac{F_1}{F_{0,x}}(\tau_m,x_t(\tau_m;0))
 +O(\varepsilon^2).
\]
If a different matched observable selects
$\tau_m(\varepsilon)=\tau_m(0)+c_t\varepsilon+\cdots$, the numerator is
$F_1+F_{0,\tau}c_t$. Ratios to a varying charge radius retain that
radius's derivative as well: for $Q=2GE/r_s$,
$(x_t/Q)_1=x_{t,1}/Q_0-x_{t,0}Q_1/Q_0^2$.
Thus the clock and radius normalization must be the same comparison
before a width coefficient is read.

These regular expansions do not make every nearby root radius
quadratic in $\sigma$. Close to the birth,
\[
 x_\pm-x_b
 =\pm\left[-\frac{2F_\tau}{F_{xx}}(\tau-\tau_b)\right]^{1/2}
   +O(\tau-\tau_b).
\]
At the unshifted birth clock, an $O(\varepsilon)$ displacement of birth
can therefore create an $O(\sqrt{\varepsilon})=O(\sigma/r_s)$ radius
split, or leave no root. A first saved root at cadence $\Delta\tau$
instead has a clock uncertainty up to $\Delta\tau$ and, at a
nondegenerate birth, a radius displacement of order
$\sqrt{\Delta\tau}$. Its record index changes discontinuously as the
width moves the birth across a checkpoint. First-record differences
are consequently not the coefficients above without event location,
matched clocks and independent control of the other errors.

\paragraph{The flat-slice clock for release from coordinate rest.}
This comparator distinguishes release from coordinate rest from the
marginally bound cloud already moving with the slice normals. On the
fixed exterior of charge $E$, set $r_s=2GE$ and write
\[
 ds^2=-dt^2+(dR+\sqrt{r_s/R}\,dt)^2+R^2d\Omega^2.
\]
The coordinate $t$ is the inward flat-slice clock, normalized at infinity.
It is neither the particle's proper time nor Schwarzschild time.
A radial timelike geodesic released with $dR/dt=0$ at
$R_0=b r_s$, $b>1$, has conserved specific energy
$e=\sqrt{1-1/b}$. With proper time $\lambda$, define
\[
 w=\sqrt{r_s/R},\qquad s=\sqrt{r_s/R-1/b},\qquad f=1-r_s/R .
\]
The normalization and energy integrals give
\[
 \frac{dR}{d\lambda}=-s,\qquad
 \frac{dt}{d\lambda}
 =\frac{e-ws}{f}
 =\frac{1+s^2}{e+ws}.
\]
The second expression is regular at the horizon. It gives the positive
fall-clock integral
\[
 \frac{\Delta t_{\rm ff}}{r_s}
 =\int_1^b
   \frac{1+(1/x-1/b)}
   {\sqrt{1/x-1/b}\,[\,e+x^{-1/2}\sqrt{1/x-1/b}\,]}\,dx .
\]
One can integrate it without assigning a finite value to a divergent
Schwarzschild-clock term. Set $x=b\cos^2\eta$ and combine that clock's
logarithm with the flat-slice transformation
$t=t_s+r_s[2\sqrt{x}+\ln((\sqrt{x}-1)/(\sqrt{x}+1))]$.
At $x=1$ the combined logarithm is finite, giving
\[
 \boxed{\frac{\Delta t_{\rm ff}}{r_s}
 =(b+2)\sqrt{b-1}\arccos(b^{-1/2})
   +(\sqrt b-1)^2+2\ln(1+b^{-1/2}) .}
\]
In particular, release from $R_0=4r_s$ gives
\[
 \Delta t_{\rm ff}
 =\bigl[2\pi\sqrt3+1+2\ln(3/2)\bigr]r_s
 \simeq12.69\,r_s .
\]
The task's $12.6r_s$ target is not the one-decimal rounding of this
expression. The marginally bound $e=1$ control is a different release:
it crosses $4r_s$ with inward coordinate speed $1/2$ and reaches $r_s$
after $(2/3)(4^{3/2}-1)r_s=14r_s/3$.

A nominal run radius is not necessarily $2GE$. If
$r_s^{\rm nom}$ labels the preparation, $R_0=4r_s^{\rm nom}$ and
$\rho=2GE/r_s^{\rm nom}$, the value on that nominal clock is
$\rho\,{\cal F}(4/\rho)$, with ${\cal F}$ the boxed function. A quoted
clock with $\rho\ne1$ is not a disagreement with the $b=4$ result.
For a pressureless boundary matched to this exterior the same geodesic
control requires no shell crossing or exchange changing the enclosed
charge. Its arrival at the charge radius is not the first trapping
birth of a centrally concentrated cloud, which can occur inside it.

\paragraph{What exclusion bounds during formation.}
Take $D$ distinct outgoing fermionic excitation channels with unit-flux
normalization on one asymptotic arrival clock. Their occupation
covariance is a contraction. In an energy band
$B=[\epsilon_1,\epsilon_2]$, $\epsilon_1\ge0$, an arrival interval
of length $\Delta t$ has density
$\Delta t\,d\epsilon/(2\pi\hbar)$ per channel. Applying
$0\leq C\leq I$ to the positive band-and-arrival number and energy
effects gives
\[
 N_{{\rm out},B}\leq\frac{D\Delta t(\epsilon_2-\epsilon_1)}{2\pi\hbar},
 \qquad
 \boxed{E_B\leq
 \frac{D\Delta t(\epsilon_2^2-\epsilon_1^2)}{4\pi\hbar}.}
\]
For example, in energy representation the positive energy effect is
$h_B^{1/2}P_{\Delta t}h_B^{1/2}$, with $h_B$ multiplication by
$\epsilon\mathbf1_B(\epsilon)$ and $P_{\Delta t}$ the Fourier arrival-time
projector. Its trace is the right-hand side; contraction bounds its
expectation by that trace. This retains frequency coherence and does
not require a thermal state, Gaussianity or independent occupancies.
It is a positive outgoing allocation, not an arbitrary signed collar
current or a local work term.

In units $c=\hbar=1$, let $\Delta t=\beta r_s$,
$q_i=\epsilon_i r_s$, $\chi=a/r_s$ and $g=G/a^2$. Since
$E=r_s/(2G)$,
\[
 \frac{E_B}{E}
 \leq\frac{D\beta g}{2\pi}(q_2^2-q_1^2)\chi^2,\qquad
 N_{{\rm out},B}\leq\frac{D\beta}{2\pi}(q_2-q_1).
\]
For fixed $D,\beta,q_1,q_2,g$, the band's outgoing energy is therefore
$O(\chi^2)$ of the charge, while its excitation count is bounded
independently of $\chi$. With $X_i=\epsilon_i/T_H=4\pi q_i$, the
energy fraction is equivalently
$D\beta g(X_2^2-X_1^2)\chi^2/(32\pi^3)$.
A logarithmically growing interval changes the bound to
$O(\chi^2\ln(1/\chi))$ under the same fixed-band and channel conditions.

To turn this into a bound on all escaping formation radiation requires
an independently controlled remainder on the same positive allocation:
\[
 \frac{E_{\rm form,out}}{E}
 \leq \frac{D\beta g}{2\pi}(q_2^2-q_1^2)\chi^2
       +\frac{E_{\rm rest}}{E}.
\]
No bound on $E_{\rm rest}$ is supplied by exclusion alone. Hard
frequencies $q\sim\chi^{-1}$, a channel count growing with
$\chi^{-1}$, or later arrivals can remove the stated suppression.
Trapped partners need not enter the outgoing allocation at all.
Consequently the capacity inequality is not a bound on the total
number of pairs formed relative to an interior reference, their
entropy, or the full light-field back-reaction. Those stronger
statements require a source-specific production and tail estimate.

\paragraph{Comparison with the supplied records [S].}
The named v1 ladder inputs define formation by the first saved
zero-holding record. Their current corrected texts,
\texttt{MASS\_ladder\_v1.1.md} and
\texttt{SCALING\_ladder\_v1.2.md}, distinguish that record's outer root
from the birth and label zeros of interpolated stored minima as
diagnostics, not event locations. This is the distinction required by
the fold calculation. They do not supply the derivatives of the
coupled $F$ needed to evaluate $b_t,b_r$, or a uniform simple-root
margin for a matched-clock expansion. The updated size source's
continuation of its width-eight member changes neither its formation
record nor its fits.

The cold reading's event-location caution agrees with the derivation;
its descriptive width fits do not establish a logarithmic width law
in place of the regular $\sigma^2$ expansion. Two certified widths at
a given size cannot test an intercept-plus-width coefficient, and
the first-record quantization remains even when each saved field is
certified. The global certification and geometric validity of the
records are separate from this conditional analytic result.

Three comparator clocks must also remain distinct. The paper's
$13.22r_s^{\rm nom}$ uses release at $4r_s^{\rm nom}$ and a charge
radius $59.05a$ with $r_s^{\rm nom}=64a$, rather than $R_0=4(2GE)$;
the conversion above specifies the comparison. The mass source
separately reports a fold at $10.8157r_s^{\rm nom}$ and first saved
root at $10.8750r_s^{\rm nom}$ for pressureless evolution of the
production's own initial profile. These are supplied readouts, not
numbers independently reproduced here. Neither is the geodesic
arrival clock $12.69(2GE)$ just derived. In particular, the cold
reading's use of the exterior boundary-arrival clock as a universal
first-formation intercept is not justified by that geodesic argument.
The exact result provides no universal pressure-correction coefficient
for the first interior birth.

The heavy-field mass and size ladders also supply no estimate of
$E_{\rm rest}$ for all light-field production. They therefore neither
extend the finite-band capacity bound to all formed pairs nor
calibrate a total light-radiation coefficient. The
finite-band statement and its explicit channel, band and arrival
conditions are the supported use.

\paragraph{Input state and comparison boundary.}
The named theorem delivery is absent in both Lane delivery folders.
The own-lane status names the served task-six(a) controls. The author
has already read parts of the v1 ladder and cold-reading inputs for
the separate item-six fallback check; no blindness to that material
is claimed. The derivations above were completed from the root
equations, the exterior metric and the covariance capacity before
any new ladder intake for this task. The source comparison above was
then made against the named v1 inputs, the corrected ladder texts and
the cold-reading event and clock clauses. It does not fit the records.
The sole independent review reproduces the three derivations; its
required fourth-moment hypothesis is included above. This note is
frozen for delivery. No manuscript statement or ladder number is
newly cleared.

\paragraph{SOURCES.}
\path{<internal>/GOAL_lane.md},
task six(a);
\path{<internal>/STATUS.md},
the named formation, clock and capacity controls;
\path{<internal>/requirements_new.txt},
the five declarations, averaging hierarchy and comparison domains;
\path{<internal>/05_star.tex}
at \texttt{a4381f19cf}, the marginally bound dust control and
coordinate-rest comparator statements;
\path{<internal>/CONDENSE_analytic_checks_v1.md},
the independently derived, frozen finite-band capacity and its scope.
The comparison uses
\path{<internal>/MASS_ladder_v1.md},
the formation definitions and section five;
\path{<internal>/SCALING_ladder_v1.md},
the formation definitions and sections 5.1 to 5.3;
\path{<internal>/MASS_ladder_v1.1.md},
the changes, definitions, formation table and section 5.1;
\path{<internal>/SCALING_ladder_v1.2.md},
the changes, definitions and sections 5.1 to 5.3;
\path{<internal>/ReviewA_reading_LADDERS_v1.md},
the event-location and comparator-clock clauses, read as a supplied
baseline, not a clearance.
The sole independent review is
\path{<internal>/redteam_SCALING_theorems_v1.md};
its kernel condition is incorporated.
No numerical ladder result or foreign proof supplies the new root
or geodesic calculation, and no manuscript text is edited.

\end{document}
